\documentclass[10pt,a4paper]{article}
\usepackage[margin=22mm]{geometry}
\usepackage[T1]{fontenc}
\usepackage{lmodern}
\ifdefined\pdfmapfile\pdfmapfile{+lm.map}\fi
\usepackage{amsmath,booktabs,graphicx}
\usepackage[expansion=false]{microtype}
\usepackage[hidelinks]{hyperref}
\renewcommand{\topfraction}{.88}
\renewcommand{\textfraction}{.10}
\renewcommand{\floatpagefraction}{.66}
% Generated by benchmarks/gears/paper.py.
\graphicspath{{\detokenize{../results/gears/figures/}}{\detokenize{../results/gears/bearing/figures/}}}
\newcommand{\SectionGain}{5.49}
\newcommand{\SectionDrift}{\ensuremath{8.92\times10^{-11}}}
\newcommand{\DutyCad}{5.87}
\newcommand{\DutyMesh}{27.75}
\newcommand{\DutyGain}{4.72}
\newcommand{\DutyGap}{32.27}
\newcommand{\DutyAxial}{1727}
\newcommand{\BearingCad}{0.139}
\newcommand{\BearingMesh}{3.780}
\newcommand{\BearingGain}{27.18}
\newcommand{\BearingDrift}{0.754}
\newcommand{\BearingMax}{3.05}
\newcommand{\BearingPeak}{66.9}
\newcommand{\BearingGuard}{44.04}
\newcommand{\MainRows}{%
Spur 24/48 & 0.299 & 2.087 & 6.99 & 2.421 & 57.1 & 1.326 \\
Helical 25/45 & 1.168 & 5.473 & 4.68 & 0.859 & 3.5 & 0.299 \\
NASA geometry 28/28 & 0.373 & 1.412 & 3.79 & 1.284 & 9.3 & 0.336 \\
Two-stage 4:1 & 1.804 & 7.582 & 4.20 & 2.117 & 32.0 & 1.854 \\
}
\newcommand{\DenseRows}{%
Helical & 3.23 & 0.825 & 0.008 & 0.9 & 0.1 \\
Reducer stage 1 & 3.76 & 1.747 & 0.004 & 28.3 & 0.1 \\
Reducer stage 2 &  & 0.899 & 0.011 & 3.2 & 0.1 \\
}

\title{\textbf{Industrial STEP Benchmarks for Implicit HHT-$\alpha$ Contact}\\
Gear Drives and a Cylindrical Roller Bearing}
\author{NurbsCONTACT research supplement}
\date{}

\begin{document}
\maketitle

\begin{abstract}
We compare native CAD and triangle-BVH contact on four loaded transmission models and a cylindrical roller bearing. Each backend uses the same time integrator, load history, body coordinates, and contact law. The test set contains a spur pair, a helical pair, a pair reconstructed from published gear-rig dimensions, a two-stage reducer, and a 14-roller bearing. The gear study measures transmission error and tooth loads through mesh and axial quadrature refinement. The bearing study compares the shaft orbit, load distribution, and a geometry-based radial clearance bound. CAD section extraction preserves the original STEP splines and reduces the cost of repeated contact probes. The delivered artifacts include the STEP solids, accepted trajectories, query counters, figures and multi-view animations.
\end{abstract}

\section{Transmission fixtures and geometry}

\begin{figure}[!htbp]
\centering
\includegraphics[width=.86\linewidth]{gear-models.pdf}
\caption{Four gear fixtures rendered from their STEP solids and accepted HHT states. Orange dots project the strongest sampled contact station of each mesh onto the camera.}
\label{fig:gear-models}
\end{figure}

The spur transmission uses 24 and 48 teeth, a 2.5 mm module, a 24 mm face, an 1800 rpm input, and 56 N\,m nominal output braking torque. The helical pair uses 25 and 45 teeth, a \(25^\circ\) helix, a 2.5 mm normal module, a 28 mm face, 1500 rpm, and 72 N\,m. The compound reducer connects a spur and helical stage through a common shaft. It uses 24/48 and 22/44 teeth, a 4:1 ratio, a \(22^\circ\) second-stage helix, an 1800 rpm input, and 148 N\,m at the output.

The fourth fixture uses the dimensions in Table I of the NASA gear-noise-rig study. It has 28 teeth per gear, a 3.175 mm module, a 6.35 mm face, and a \(20^\circ\) pressure angle. We included the reported 23 and 25 \(\mu\)m tip relief beginning at a \(24^\circ\) roll angle. We use 2000 rpm and a nominal torque of 71.77 N\,m. The bench dimensions and relief inform the generated STEP, while the stiffness, shaft inertia, and excitation are specified by this benchmark. The comparison uses numerical references under this model. \href{https://ntrs.nasa.gov/citations/19910013253}{Oswald et al., NASA gear-noise-rig report}.

Each gear is a closed STEP solid with an axial bore, repeated active flanks, tooth tips, roots, and end faces. The working faces are cubic spline interpolants of transverse involutes with axial lead and optional relief. The probe reuses an original working flank through rotational tooth symmetry. A sampled validation compares the round-tripped STEP flank against the analytical generating curve. The per-flank fit stays within 2 \(\mu\)m over that validation grid. The roots use simplified transitions and do not supply bending stresses.

\section{Shared contact and integration}

For each stage let \(s\) be distance along the transverse line of action, \(r_b\) the driver base radius, \(h\) its twist per unit face width, and \(p_b\) the base pitch. At axial station \(z\), tooth pair \(j\) is sampled at the nominal position
\begin{equation}
s_j = r_b(\omega t+h z)+j p_b.
\end{equation}
The two flanks undergo their actual rigid transforms from the HHT states. At the candidate station we solve two short ray intersections on the trimmed working flanks. The signed gap is \((t_B-t_A)\cos\beta_b\), where \(\beta_b\) is the base helix angle. Both pipelines use the same finite-path and endpoint tapers, contact stations, and linear compliance of \(1.6\times10^8\) N/m per fully weighted tooth. This is a projected normal contact model on the driving flanks. The projected station gaps are evaluated independently on native CAD during postprocessing.

The native queries intersect fixed-axial STEP isocurves. We extract their nonrational cubic segments and evaluate them by Horner's method in a restricted Newton solve. The triangle backend uses fresh OCCT tessellations of the same trimmed faces and two-sided finite ray intersections with a median-split AABB hierarchy. On the helical ablation, extracting the exact section basis gives a \SectionGain-fold speedup relative to repeated OCCT surface evaluation. The largest transmission difference was \SectionDrift{} \(\mu\)m.

The common HHT-\(\alpha\) solver uses \(\alpha=-0.1\), shaft-angle deviations from nominal rotation, and an implicit force residual. Shaft inertias use annular estimates. The load history contains a rise near 0.10 s, a release near 0.24 s, and a superposed periodic component. The driver is coupled to its nominal angle by a torsional spring and damper. Positive load maintains driving-flank engagement. The models resolve effective normal and axial loads. Friction, rolling lubrication, tooth fatigue, and housing flexibility are outside the reduced equations.

\section{Measured results}

\begin{table}[!htbp]
\small\centering
\caption{0.60 s gear runs at 20 \(\mu\)s. Medians use three repetitions. The mesh deflection is 1.5625 \(\mu\)m. TE and force columns compare the mesh with matched-step CAD. The last column compares CAD at 20 and 10 \(\mu\)s. The reducer entries use the larger stage error per metric.}
\label{tab:primary-gears}
\begin{tabular}{lrrrrrr}
\toprule
Fixture & CAD (s) & Mesh (s) & Ratio & TE (\(\mu\)m) & Force (\%) & CAD step TE (\(\mu\)m)\\
\midrule
\MainRows
\bottomrule
\end{tabular}
\end{table}

\begin{figure}[!htbp]
\centering
\includegraphics[width=.94\linewidth]{gear-cost-accuracy.pdf}
\caption{Runtime ratios and transmission RMS differences over the three tessellation settings. All settings use the same contact projection and numerical integration.}
\end{figure}

The primary matrix uses one axial station for spur gears and seven stations for helical gears. Table~\ref{tab:primary-gears} separates temporal discrepancy from the matched-step effect of changing the flank representation. It measures transmission error in base-circle displacement units. The force metric is the RMS difference divided by the mean CAD force after 20 ms. Mesh fineness changes both the cost per ray and the number of implicit residual evaluations. A small RMS transmission discrepancy can accompany substantially different force peaks.

\begin{table}[!htbp]
\small\centering
\caption{0.36 s axial refinement. Runtime ratios compare CAD and the fine mesh at 31 axial stations. Error columns compare the fine mesh with CAD at 31 stations and CAD at 31 stations with CAD at 63, respectively. Spur stages use a single station.}
\label{tab:axial}
\begin{tabular}{lrrrrr}
\toprule
Fixture & Ratio & Mesh TE (\(\mu\)m) & Axial TE (\(\mu\)m) & Mesh force (\%) & Axial force (\%)\\
\midrule
\DenseRows
\bottomrule
\end{tabular}
\end{table}

With helical contact, the axial station count changes the force waveform. We measured additional 31- and 63-station trajectories to quantify the sensitivity independently of surface faceting. The fine mesh and native CAD use the same 31 stations in this comparison. Figure~\ref{fig:axial} plots the corresponding force and transmission histories, and Table~\ref{tab:axial} retains the residual sensitivity to axial quadrature.

\begin{figure}[!htbp]
\centering
\includegraphics[width=.86\linewidth]{gear-axial-refinement.pdf}
\caption{Refinement at matched contact stations and at higher native axial sampling. The same projected contact model is used in both refinement directions.}
\label{fig:axial}
\end{figure}

The 3 s reducer trajectories complete 150,000 nominal steps, covering 2160 stage-one and 990 stage-two tooth-passing periods. Integration took \DutyCad{} s on CAD and \DutyMesh{} s on the fine mesh, a factor of \DutyGain. The smallest native sampled clearance was \DutyGap{} \(\mu\)m on the CAD trajectory. The largest native axial load on the helical stage was \DutyAxial{} N. The gear clearance evidence applies to the selected working-flank probes and recorded times.

\begin{figure}[!htbp]
\centering
\includegraphics[width=.88\linewidth]{reducer-duty.pdf}
\caption{Long reducer runs. Filled bands give the observed range in 20 ms windows. Force traces mark the per-window mean and gap traces mark the per-window minimum.}
\end{figure}

\section{Cylindrical roller bearing}

The bearing uses the 45 mm bore, 85 mm outer diameter, and 19 mm width of the NU 209 envelope. The roller count, race radii, and stiffness are defined here. Fourteen 8 mm rollers share a prescribed 787.5 rpm cage driven by an 1800 rpm shaft. The HHT state contains the two radial shaft translations. A massless radial series-compliance model balances the opposing roller-to-race loads. A radial force step and rotating imbalance drive the shaft. \href{https://www.emarketplace.in.skf.com/cylindrical-roller-bearing/nu-209-ecp}{SKF catalog dimensions}.

The native method uses analytic radii extracted from the three round STEP surfaces. The triangle method intersects equivalent radial probes through meshes of those faces. For a linearly reconstructed displacement \(q\) of the shaft center, the total directional room remaining after the roller diameter is at least
\[
R_o-R_i-2r-\max(\|q_{\rm start}\|,\|q_{\rm end}\|).
\]
Both variants use this interval bound, evaluated with outward rounding. The bound describes the directional radial model under prescribed cage kinematics.

In the 0.36 s test the CAD integration median was \BearingCad{} s and the fine-mesh median was \BearingMesh{} s. The speed ratio was \BearingGain. The RMS shaft-orbit difference was \BearingDrift{} \(\mu\)m, with a maximum of \BearingMax{} \(\mu\)m. The RMS difference in maximum roller force was \BearingPeak{} N. The smallest common radial interval bound across both configurations was \BearingGuard{} \(\mu\)m.

\begin{figure}[!htbp]
\centering
\includegraphics[width=.87\linewidth]{bearing-loads.pdf}
\caption{STEP bearing assembly and load zone, with matched-step shaft orbit and roller load traces.}
\end{figure}

\IfFileExists{gears-experiment-results.tex}{%
  % Generated by benchmarks/gears/experimental/paper_section.py.
\newcommand{\GearExperimentPath}{\detokenize{../results/gears/experimental/figures/}}
\newcommand{\NasaK}{60}
\newcommand{\NasaZeta}{0.20}
\newcommand{\UnswK}{50}
\newcommand{\UnswZeta}{0.20}
\newcommand{\UnswShift}{0.500}
\newcommand{\NasaStepDelta}{0.28}
\newcommand{\UnswStepDelta}{0.009}
\newcommand{\NasaInertiaSensitivity}{76.4}
\newcommand{\NasaInertiaRMSE}{164.7}
\newcommand{\NasaLowLoadFloor}{247}
\newcommand{\UnswMeasuredSecond}{6.08}
\newcommand{\UnswNativeSecond}{0.37}
\newcommand{\NasaExperimentRows}{%
7(a) & 2000 & 33.73 & 909 & 63.1 & 66.7 & Fit \\
7(b) & 4000 & 11.48 & 462 & 146.7 & 150.0 & Held out \\
7(c) & 6000 & 33.73 & 1078 & 148.8 & 151.7 & Held out \\
7(d) & 6000 & 67.46 & 1785 & 90.8 & 96.4 & Held out \\
}
\newcommand{\UnswExperimentRows}{%
297.8 & 10.36 & 2.74 & 1.95 & 1.95 & Fit \\
298.9 & 5.38 & 1.92 & 1.46 & 1.45 & Held out \\
596.2 & 15.32 & 5.17 & 4.35 & 4.39 & Held out \\
}
%
  \section{External gear-rig comparisons}

The simulated gear benchmarks above control numerical resolution and the contact representation. We additionally test the reduced contact dynamics against a single-tooth force measurement and against two-shaft encoder measurements. Stiffness and damping are selected from a fixed grid using one operating condition per rig. The mesh and B-rep runs share the selected parameters. Scoring on held-out conditions uses the same parameters and phase conventions.

\subsection{NASA single-tooth dynamic force}

We use the reconstructed 28/28-tooth CAD fixture described above and constant torque at the values indicated in Fig.~7 of NASA TM-103232. Four torsional coordinates represent the motor, pinion, wheel, and load. The model uses the inertia values printed in SI units in Table II and the reported 341 and 12,700 N\,m/rad input and output torsional stiffnesses. An equivalent normal compliance represents the working contact. The gearbox stiffness entry from the reference is distinct from the effective normal tooth stiffness fitted here. The shaft damping is set to 0.05 of critical for the adjacent-inertia pair. The selected effective tooth stiffness is \NasaK{} MN/m with damping ratio \NasaZeta{}. Both values lie on the prescribed NASA calibration grid boundary.

We digitize the four experimental traces, using only the solid force curves. The points trace the published plot with an ordinate resolution of roughly 6 N per source pixel. Thick and overlapping strokes add local reading uncertainty. The comparison uses the driven tooth roll angle from the actual native or meshed witness. We use the simulated single-pair forces, average complete simulated engagements in the last half of each run, and hold the roll-angle origin fixed. The scored range ends at \(29.5^\circ\), with the lower bound at \(12^\circ\) for (a)--(c) and \(12.625^\circ\) for (d). The restricted endpoint of (d) follows the identifiable trace of the plotted measurement. \href{https://ntrs.nasa.gov/citations/19910013253}{NASA TM-103232}.

\begin{table}[!htbp]
\centering\small
\caption{Single-tooth normal force relative to the digitized experiments. Peak refers to the scored roll-angle range. RMSE is computed on the common \(0.125^\circ\) grid.}
\label{tab:nasa-exp}
\begin{tabular}{lrrrrrl}
\toprule
Panel & rpm & Torque (N\,m) & Measured peak (N) & B-rep RMSE (N) & Mesh RMSE (N) & Role\\
\midrule
\NasaExperimentRows
\bottomrule
\end{tabular}
\end{table}

\begin{figure}[!htbp]
\centering
\includegraphics[width=.88\linewidth]{\GearExperimentPath nasa-single-tooth.pdf}
\caption{NASA single-tooth force. Dots digitize the experimental trace. Solid and dashed lines are the HHT runs using native STEP faces and triangles, respectively.}
\label{fig:nasa-exp}
\end{figure}

The low-load 4000 rpm measurement shows loss of tooth contact. Within the central \(19.5^\circ\) to \(21.5^\circ\) interval the native single-tooth force remains above \NasaLowLoadFloor{} N, leaving the loss of contact unpredicted. In the 6000 rpm, 94\% torque case, halving the native time step changes the force profile by \NasaStepDelta{} N RMS. The paired SI and customary inertia entries in NASA Table II imply a tenfold difference when the customary entries are interpreted as force-inch units. A separate sensitivity run multiplies all four inertias by ten with the other settings fixed. At 6000 rpm and 47\% torque it changes the force profile by \NasaInertiaSensitivity{} N RMS, giving an error of \NasaInertiaRMSE{} N against the digitized measurement. The main comparison retains the published SI entries.

\IfFileExists{nasa-revised-results.tex}{%
  % Generated by benchmarks/gears/experimental/nasa_revision_artifacts.py.
\newcommand{\NasaRevisedPath}{\detokenize{../results/gears/experimental/revised/static-matched/figures/}}
\newcommand{\NasaRevisedStiffness}{67.5}
\newcommand{\NasaRevisedDamping}{0.10}
\newcommand{\NasaRevisedKBeam}{172.8}
\newcommand{\NasaRevisedFrequencies}{6.73, 665.63, 4106.59}
\newcommand{\NasaRevisedGearboxFrequencies}{6.56, 524.90, 1220.79}
\newcommand{\NasaRevisedStaticSpan}{15.87, 24.88}
\newcommand{\NasaRevisedStaticTarget}{16.00, 25.00}
\newcommand{\NasaRevisedLowMin}{125.9}
\newcommand{\NasaRevisedLowOverlap}{0}
\newcommand{\NasaRevisedObservedSpan}{19.25, 22.00}
\newcommand{\NasaBeamInternalSpan}{23.50, 24.38}
\newcommand{\NasaRevisedSmallestGap}{29.33}
\newcommand{\NasaRevisedStepB}{0.42}
\newcommand{\NasaRevisedStepD}{0.16}
\newcommand{\NasaQuarterStepDelta}{1.33}
\newcommand{\NasaQuarterStepMin}{125.21}
\newcommand{\NasaQuarterStepSamples}{51,409}
\newcommand{\NasaRevisedTimeBrep}{1.53}
\newcommand{\NasaRevisedTimeBvh}{6.94}
\newcommand{\NasaRevisedSpeed}{4.54}
\newcommand{\NasaRevisedDualShares}{10.3, 29.9, 60.7}
\newcommand{\NasaRevisedRows}{%
7(a) & 63.1 & 67.8 & 125.4 & 63.2 & 63.9 & -3.5 \\
7(b) & 146.7 & 173.7 & 226.4 & 203.2 & 201.9 & +0.6 \\
7(c) & 148.8 & 164.7 & 332.6 & 174.0 & 176.0 & -17.6 \\
7(d) & 90.8 & 90.7 & 260.1 & 75.7 & 81.6 & -4.7 \\
}
%
  \subsection{Reconstructed gear compliance and the loss-of-contact test}

We rebuilt both STEP flanks from the average shape of the four traces per gear in Fig.~3 of NASA TM-103232. The plotted profiles provide the relative shape. Table~I fixes the tip amplitudes at 23 and \(25\,\mu\mathrm m\). This reconstruction shares the same geometric profile across all teeth of a gear. The native and triangle runs use the same exported STEP surfaces and trimmed line-of-action stations. We converted the customary inertia column to SI using \(1\,\mathrm{lbf\,in\,s^2}=0.112985\,\mathrm{kg\,m^2}\). The \href{https://ntrs.nasa.gov/citations/19930018861}{DANST manual} independently lists the customary magnitudes. The equivalent gearbox stiffness in Table~II includes shaft flexibility. It enters the modal audit separately from the tooth pair compliance.

At each station we compute the compliance of the two gear teeth as tapered beams with bending, shear, and radial compression, and add a line-contact compliance. Let \(g_i\) denote the CAD gap, \(g_0\) the fixed gap offset, \(C_i\) the sum of those compliances, and \(C_b\) a constant series compliance accounting for deformation outside the reduced tooth model. With geometric endpoint weight \(w_i\), rolling normal force satisfies
\[
 k_i=(C_i+C_b)^{-1},\qquad
 F_i=w_i\max\!\left(0,k_i(g_0-g_i)-c_i\,\dot g_i\right)
 \quad\text{for }g_i<g_0 ,
\]
and vanishes for \(g_i\ge g_0\). The damping coefficient uses a ratio of \NasaRevisedDamping{} and the reduced inertia. The gap derivative includes the advance of the nominal gear angle and the generalized velocities. We obtain its signed lever arms from the tangent of each native isocurve or the normal of the intersected triangle. The load tangent for the HHT solve uses the same gap gradient. The query keeps the finite-flank trim and the original \(2\,\mu\mathrm m\) sampled working-flank guard.

The independent beam and line-contact estimate gives \NasaRevisedKBeam{} MN/m at the pitch point. For the series-compliance variant we solve torque balance over 224 mesh phases at each load. We select one stiffness using only the approximate static DANST plateau in Fig.~7(c), near roll angles \(\NasaRevisedStaticTarget\) degrees. The selected pitch stiffness is \NasaRevisedStiffness{} MN/m, with plateau endpoints \(\NasaRevisedStaticSpan\) degrees. The reference plateau is a published calculation. Its overlapping strokes limit the precision of the endpoints. The static torque balance residual stays below \(10^{-6}\,\mathrm{N\,m}\). For the native runs, the fractions of phases with two loaded tooth pairs are \NasaRevisedDualShares{} percent at 11.48, 33.73, and 67.46 N\,m, respectively.

The four dynamic comparisons use the same stiffness and damping for both geometric backends. The gearbox-stiffness linearization with the converted inertias gives \(\NasaRevisedGearboxFrequencies\) Hz. The report lists 6.56, 52.5, and 1220 Hz. The middle entry is inconsistent with our four-inertia reconstruction. The single loaded pair at the calibrated pitch stiffness gives \(\NasaRevisedFrequencies\) Hz. Table~II lists the highest DANST frequency at 5821 Hz. A sensitivity run imposed that value on the frozen single-pair linearization. It produced an early zero-force interval with zero overlap with the measured zero bins.

\begin{figure}[!htbp]
\centering
\includegraphics[width=.87\linewidth]{\NasaRevisedPath nasa-revised-force.pdf}
\caption{NASA single-tooth force with converted inertias, digitized profile shapes, and the static-calibrated series compliance. The shaded region in 7(b) marks the measured zero-force interval on the scoring grid.}
\label{fig:nasa-revised}
\end{figure}

\begin{table}[!htbp]
\centering\small
\caption{Waveform RMSE in newtons for the successive NASA model changes. Peak error refers to the static-calibrated native model.}
\label{tab:nasa-revised}
\resizebox{\linewidth}{!}{%
\begin{tabular}{lrrrrrr}
\toprule
Panel & Original B-rep & Inertias only & Beam B-rep & Calibrated B-rep & Calibrated BVH & B-rep peak err. (\%)\\
\midrule
\NasaRevisedRows
\bottomrule
\end{tabular}}
\end{table}

In 7(b) the digitized measured zero range is \(\NasaRevisedObservedSpan\) degrees. The pure beam model predicts an internal zero-force interval at \(\NasaBeamInternalSpan\) degrees. The calibrated B-rep result overlaps \NasaRevisedLowOverlap{} percent of the measured zero bins and reaches a minimum of \NasaRevisedLowMin{} N on them. The observed phase discrepancy persists under step refinement. Halving the step changes the calibrated B-rep waveform by \NasaRevisedStepB{} N RMS in 7(b) and \NasaRevisedStepD{} N RMS in 7(d). The calibrated runs retain a smallest recorded gap of \NasaRevisedSmallestGap{} micrometers. The remaining force errors motivate an independently identified dynamic model of the shafts and supports.

Refining 7(b) again to \(0.25\,\mu\mathrm s\) changes its waveform by \NasaQuarterStepDelta{} N RMS relative to the nominal step. Before phase averaging, \NasaQuarterStepSamples{} saved tooth samples occupy the measured zero-force interval. Their minimum force is \NasaQuarterStepMin{} N. The output cadence in this audit is \(1\,\mu\mathrm s\).

In three alternating, serial timing repeats of 7(b), each simulating 0.16 s, median HHT integration takes \NasaRevisedTimeBrep{} s for the native surfaces and \NasaRevisedTimeBvh{} s for the triangle BVH, a factor of \NasaRevisedSpeed{}. Timings exclude geometry setup and output.
%
}{}

\subsection{UNSW encoder transmission error}

The UNSW test gear pair has 27/44 teeth, a 2 mm module, 5 mm face width, \(20^\circ\) pressure angle, and 10 mm bore radii. We use three healthy records from the original 200 kHz dataset. The input reference phase is set by the tachometer. The 3600-pulse encoder signals on both shafts are converted to continuously interpolated angular measurements using linearly interpolated zero crossings. The base-radius-scaled difference gives the transmission error. We remove the block mean and average complete hunting-tooth periods of 44 input revolutions, preserving the mesh-synchronous waveform. Repeating the extraction with falling edges changes the three mean waveforms by less than \(0.012\,\mu\)m RMS. The reported block range shows variation within each recording.

The reconstructed solver uses two effective shaft inertias and ideal involute STEP geometry. Encoder transfer dynamics and as-manufactured deviations are omitted from this model. The effective stiffness \UnswK{} MN/m, damping ratio \UnswZeta{}, and the phase offset of \UnswShift{} mesh cycles are set using the nominal 5 Hz and 10 N\,m record. The two held-out records use measured speeds and mean input torques. Halving the native step in the faster held-out case changes the mesh-synchronous waveform by \UnswStepDelta{} \(\mu\)m RMS. \href{https://doi.org/10.26190/UNSWORKS/24007}{Test geometry in Chin's thesis}. \href{https://doi.org/10.17632/p72x3m92cv.1}{CC BY 4.0 dataset}.

\begin{table}[!htbp]
\centering\small
\caption{UNSW AC transmission error. RMS and RMSE are computed on 128 uniformly spaced points over one mesh cycle.}
\label{tab:unsw-exp}
\begin{tabular}{rrrrrl}
\toprule
Measured rpm & Input N\,m & Data RMS (\(\mu\)m) & B-rep RMSE (\(\mu\)m) &
Mesh RMSE (\(\mu\)m) & Role\\
\midrule
\UnswExperimentRows
\bottomrule
\end{tabular}
\end{table}

\begin{figure}[!htbp]
\centering
\includegraphics[width=.88\linewidth]{\GearExperimentPath unsw-encoder-dte.pdf}
\caption{UNSW DTE and mesh harmonics at the calibration setting and two held-out settings. The shaded band is the range across complete hunting-tooth blocks of the encoder record.}
\label{fig:unsw-exp}
\end{figure}

At the faster held-out setting the measured second mesh harmonic is \UnswMeasuredSecond{} \(\mu\)m and the native result is \UnswNativeSecond{} \(\mu\)m. The held-out errors characterize the full reduced model. Residual discrepancies can include contact-law approximations, rig dynamics, and manufacturing deviations. Replacing the triangulated faces with native STEP geometry isolates one numerical contribution within those settings.
%
}{}

\section{Reproduction}
The \path{nc_gears} and \path{nc_roller_bearing} targets generate the working geometry and run the contact models. \path{results/gears/README.md} gives the Python analysis and rendering commands. The protocol and timing files record each deflection, step size, repeat and axial sampling setting. Timed integration excludes STEP loading, contact representation setup, and disk output. Setup is recorded separately. The models and runs use the project OCCT 7.9.3 Release build.

\end{document}
